\section{Receptive fields: the anisotropy is real, structural, and invisible to the obvious metric}
\label{sec:erf}

A state space backbone reads an image as a one-dimensional sequence. Under a
row-major flattening two vertically adjacent pixels are separated in that
sequence by the width of the token grid, while two horizontally adjacent
pixels are neighbours. The effective receptive field should therefore be
elongated along the scan direction. This is asserted in the comparison
literature, including in our own earlier version, without measurement.

The claim turns out to be true, false, and true again depending on how it is
measured, and the resolution of that is the useful result.

\begin{figure}[tbp]
  \centering
  \includegraphics[width=\linewidth]{fig2_erf.pdf}
  \caption{Effective receptive fields. Columns are input conditions and the
  rightmost panel of each row shows how the two metrics converge with the
  number of images averaged. The untrained Vim-S field is visibly elongated;
  the trained one is not obviously so, and the covariance index does not
  register the difference.}
  \label{fig:erf}
\end{figure}

\begin{table}[t]
\centering
\caption{ERF anisotropy at $n=512$ images. The covariance index separates no model; the axis-wise decay ratio and the untrained condition do.}
\label{tab:e2}
\begin{tabular}{llrrr}
\toprule
Model & Condition & Anisotropy & Principal angle (deg) & Decay ratio \\
\midrule
DeiT-S & natural & 1.009 & 34.06 & 1.048 \\
DeiT-S & random\_init & 1.005 & 33.50 & 1.013 \\
CMT-S & natural & 1.028 & 80.90 & 1.023 \\
CMT-S & random\_init & 1.025 & 36.48 & 1.107 \\
Vim-S & natural & 1.045 & 14.03 & 1.345 \\
Vim-S & random\_init & 1.979 & 0.07 & 1.899 \\
\bottomrule
\end{tabular}
\end{table}


\subsection{The covariance index separates nothing}

The natural way to quantify elongation is the ratio of the principal axes of
the second-moment matrix of the field. Table~\ref{tab:e2} reports it, and on
natural images all three models sit close to unity. On this metric the three
architectures are not distinguishable, and the claim as usually stated is not
supported.

The metric is at fault rather than the claim. A second moment weights each
point by the square of its distance from the centre, so the value is
dominated by the broad and nearly isotropic skirt of the field rather than by
its core. Whatever the scan does to the shape near the peak is averaged away
by a large amount of nearly circular background.

\subsection{An axis-wise decay ratio does separate them}

Measuring instead the ratio of vertical to horizontal decay slopes in the
neighbourhood of the peak, Vim-S sits at \VimDecayNatural{} while the two
attention models sit near unity---a separation of roughly thirty percent.

That separation is an order of magnitude larger than the measurement
uncertainty. The experiment was run twice over different image sets, and
comparing the two at the same sample size the Vim-S decay ratio moves by
\VimDecayRunDrift{} percent. The gap being claimed is ten times the drift
that reproduces it.

The difference between CMT-S and DeiT-S on this metric is inside that
uncertainty, and we do not claim an ordering between them.

\subsection{The clearest evidence comes from a model that has learned nothing}

Take Vim-S with random weights and no training at all. Its receptive field
has an anisotropy index of \VimAnisoRandom{}---an axis ratio close to two to
one---and the principal axis of that ellipse lies within
\VimAngleRandom{} degrees of the horizontal. The decay ratio in the same
condition is \VimDecayRandom{}, larger than in the trained model. Both hold
across two runs with different random initialisations and at every sample
size from sixteen images to five hundred and twelve.

This is the claim measured literally, and it locates the cause. The
anisotropy is not something the network learns; it is imposed by the
row-major flattening and the exponential state decay along the sequence, and
it is present before any gradient has been applied. Training does not create
the effect. Training partially cancels it, which is why the trained decay
ratio is the smaller of the two and why the covariance index---already the
least sensitive of the three measurements---finds nothing left to see.

An untrained network is usually a control. Here it is the condition in which
the phenomenon is largest, because it is the condition in which nothing has
had a chance to compensate for it.
